\section{Evaluation Protocol}

\subsection{Primary Deployment Metrics}

The report uses the repo's downstream warm-start metrics rather than offline loss alone:
\begin{itemize}
\item solver success rate,
\item solver iterations,
\item solver solve time,
\item total DT time = warm-start inference + solve time,
\item warm-start acceptance rate,
\item fallback count,
\item projection fraction,
\item projection total magnitude,
\item projection max-step magnitude.
\end{itemize}

A central lesson of the project is that validation loss and deployment metrics must be separated. Several later checkpoints improved offline and still failed to improve solver-facing behavior.

\subsection{Benchmark Scope}

The final narrowed benchmark uses two frozen Oval gates:
\begin{itemize}
\item no-obstacle gate,
\item 1--4 obstacle gate.
\end{itemize}

Most completed controlled evaluations use IPOPT rather than FATROP. This was a practical choice: FATROP remained the main generation path for clean training data, but IPOPT was used for fixed-gate diagnostic evaluations because it was faster and more reliable for repeated interactive benchmarking.

\subsection{Important Evaluation Caveats}

Three caveats are required for faithful interpretation.

First, an early no-obstacle evaluation was confounded by repeated starts. The evaluation path was effectively repeating the same start configuration, so an initial 10-scenario no-obstacle gate was not actually testing 10 distinct start anchors. This was later corrected.

Second, the \texttt{lap\_time\_s} field stored in evaluation outputs corresponds to the optimizer objective value rather than a pure terminal-time readout, so we do not use it as a primary KPI in this report.

Third, acceptance rates depend strongly on the rejection gate. A corrected relaxed-gate diagnostic is useful for understanding rollout quality, but it does not by itself show solver benefit.
